\section{Introduction}

Predictive systems increasingly combine multiple auxiliary contexts, including
retrieved examples, historical observations, and neighboring records. The
usefulness of each context depends on what the predictor has already received.
A context that improves a prediction in isolation can become redundant after
similar information is selected, while another candidate with lower standalone
value can correct the remaining error. A selection rule that assigns each
candidate a fixed score cannot represent this change.

We model the value of a candidate directly as its marginal reduction in
end-task loss,
\begin{equation}
 m(j\mid A)=F(A\cup\{j\})-F(A),
\end{equation}
where $A$ is the current context set and $F$ measures the expected loss
reduction relative to an anchor-only predictor. The marginal depends on the
predictor's current residual and on the direction and magnitude of the
prediction change induced by the candidate. It therefore differs from a
standalone score $u(j\mid\emptyset)$: the same candidate can be useful at one
selection state and redundant or harmful at another.

This distinction is structural rather than terminological. A perfect
standalone score fixes candidate values before selection. A sequential oracle
recomputes the true marginal after each addition. Their gap,
\begin{equation}
 H_{\mathrm{state}}
 =
 \frac{G_{\mathrm{seq}}-G_{\mathrm{stand}}}{G_{\mathrm{seq}}},
\end{equation}
measures the value created solely by conditioning on the selected set. We
call this quantity structural headroom. Controlled experiments show that
redundancy enlarges it, and that a learned marginal-utility router can surpass
an exact standalone ranking when redundancy is high.

Learning to estimate marginal utility is harder than predicting a fixed
candidate score. The router must distinguish candidates whose utilities are
close, and it must know how each candidate changes the fixed predictor. As the
candidate pool grows, the utility gap between the best alternatives shrinks.
Cached context representations, in turn, omit information carried by the
predictor's response to a candidate. These two properties set the design
requirements for a practical router.

We introduce Response-Aware Marginal-Utility Routing (R-MUR). A low-cost cached
score first narrows the candidate pool. The fixed predictor is evaluated only
on this shortlist, and compact summaries of its response are used to rerank the
retained candidates. Training combines marginal-utility regression with a
pairwise objective weighted by the utility gap, which aligns estimation with
the final selection decision. The design follows directly from the two
properties above: predictor-response features expose the change that enters
the marginal, and ranking calibration targets small decision margins.

The resulting evidence chain is deliberately narrow, and it has a boundary
that we characterize rather than hide. Controlled studies isolate how
redundancy and candidate-set growth change the selection problem. Across
standard traffic benchmarks, R-MUR improves learned context selection over
cached and standalone baselines, while a multi-target study shows that
state-conditioning headroom persists across sensors. These results establish
state-conditioned marginal utility as a practical decision object for budgeted
multi-context prediction and identify predictor response as a key signal for
learning it. We then test the same framework against a subset-capable
nonlinear backbone and against a second task family. The nonlinear backbone
preserves structural headroom but shrinks the share of it that a router can
realize: response-free learned routers become indistinguishable from random
selection and R-MUR's advantage over them is no longer statistically resolved.
The second family, whose candidate demonstrations are not redundant, is already
solved by nearest-neighbour relevance. Together with the smooth-loss analysis,
these results locate the regime in which sequential marginal routing pays off:
candidates must overlap in what they contribute, and the fixed predictor must
retain a residual that the router can observe.

\paragraph{Scope and companion manuscripts.} This paper reports the method and
the generalisation evidence: the pipeline, the ten-method baseline matrix, the
strong-backbone boundary, the second task family, the compute profile and the
adaptive rule. The identifiability theory --- why the predictor's response
decides whether utility can be estimated, and why stronger predictors make
routing harder --- is developed and proved in a companion manuscript. The
summary propositions reproduced here (\cref{sec:analysis}) are stated so that
this paper's predictions can be read without it.
